% Generated draft: requires mathematical and source review.
\documentclass[11pt]{article}
\usepackage[margin=25mm]{geometry}
\usepackage{fontspec}
\setmainfont{Latin Modern Roman}
\usepackage{amsmath,amssymb,mathtools,unicode-math}
\setmathfont{Latin Modern Math}
\usepackage{enumitem,needspace,xurl,hyperref}
\setlength{\parindent}{0pt}
\setlength{\parskip}{5pt}
\setlength{\emergencystretch}{3em}
\newcommand{\N}{\mathbb{N}}
\newcommand{\Z}{\mathbb{Z}}
\newcommand{\Q}{\mathbb{Q}}
\newcommand{\R}{\mathbb{R}}
\newcommand{\C}{\mathbb{C}}
\newcommand{\F}{\mathbb{F}}
\newcommand{\A}{\mathbb{A}}
\newcommand{\PP}{\mathbb{P}}
\newcommand{\E}{\mathbb{E}}
\newcommand{\eps}{\varepsilon}
\newcommand{\dd}{\,\mathrm{d}}
\begin{document}
\section*{UnsolvedMath ID: 30000166}
\textbf{Source problem code:} catalogue entry
\subsection{Definitions and mathematical statement}
Let \(h\in\Z_{>0}\), let \(\varphi(h)=|\{1\leq a\leq h:\gcd(a,h)=1\}|\), and let \(\mu\) be the M\"obius function: \(\mu(1)=1\), \(\mu(d)=0\) if a prime square divides \(d\), and \(\mu(d)=(-1)^k\) if \(d\) is a product of \(k\) distinct primes. For \(\tau\) in the complex upper half-plane set \(q=e^{2\pi i\tau}\) and
\[
 \eta(\tau)=e^{2\pi i\tau/24}\prod_{m\geq1}(1-q^m),\qquad
 F_h(\tau)=\eta(h\tau)^{\varphi(h)}\prod_{d\mid h}\eta(d\tau)^{-\mu(d)}.
\]
The divisors \(d\) are positive divisors. Write
\[
 \alpha_h=\frac{h\varphi(h)-\sum_{d\mid h}d\mu(d)}{24},\qquad
 F_h(\tau)=e^{2\pi i\alpha_h\tau}\sum_{n\geq0}a_h(n)q^n.
\]
Equivalently the coefficients are defined by the formal power series
\[
 \sum_{n\geq0}a_h(n)q^n
 =\prod_{m\geq1}(1-q^{hm})^{\varphi(h)}
   \prod_{d\mid h}\prod_{m\geq1}(1-q^{dm})^{-\mu(d)}.
\]
All exponents of the factors are integers and each coefficient depends on only finitely many factors, so this defines \(a_h(n)\in\Z\). Factoring out the initial rational power of \(q\) makes the Fourier-coefficient convention explicit without requiring \(\alpha_h\) to be an integer.

The cleaned conjecture, attributed in the original record to K.\ Saito, asks whether
\[
 \forall h\in\Z_{>0}\ \forall n\in\Z_{\geq0},\qquad a_h(n)\geq0.
\]
The specifically emphasized case is to prove or disprove this assertion for every \(h\) divisible by at least two distinct primes. Nonnegative means zero is allowed; it is not a strict-positivity assertion. Natural numbers are interpreted here as positive because the displayed eta product is not defined at \(h=0\).

The original wording singles out that multiple-prime case. The verification field reports that the official report was checked and that Saito's conjecture and the remaining composite case were then stated explicitly. The dated literature assessment, checked 2026-08-21, instead describes an extraction containing unnamed conjectures with missing statements. That description appears to concern an earlier or different extraction and does not assess the explicit eta-product formulation on its own. No proof, classification of already established cases, or independently verified current open status is supplied by this draft.
\subsection{Short English statement}
Saito's conjecture predicts that every coefficient of the displayed eta product, after its initial power of \(q\) is separated, is nonnegative. The source particularly asks for the case where \(h\) has at least two distinct prime divisors. The explicit repaired statement is retained together with the older, apparently incompatible literature note.
\subsection{Sources}
\begin{itemize}
\item Local UnsolvedMath registry, permanent ID 30000166, source code catalogue entry, ``Positivity of Saito's Eta Products.'' The original statement attributes the eta product to K.\ Saito; the cleaned statement and verification field give the coefficient-positivity conjecture.
\item \emph{Gitter und Anwendungen (2005)}, Oberwolfach Reports, pp.\ 5--67, DOI 10.4171/owr/2005/01: \url{https://doi.org/10.4171/owr/2005/01}. This is the sole supplied literature reference, labelled ``original OWR report.''
\item Literature triage checked 2026-08-21; local record updated 2026-09-16. The triage describes unnamed missing conjectures, whereas the verification field reports a direct-source repair. No external source has been visited or independently checked here.
\end{itemize}
\end{document}
